\subsection{Odd terminals with a nonunit anchor}\label{loc:S32}
Let $2\le a<b_1<c_1<\cdots<b_t<c_t$, $t\ge1$, $b_1\ge4$.
Use the actual inward anchor $\bar f_{a}$ and
\[
 P=\begin{cases}\prod_{i} h_{b_i}f_{c_i},&b_1\ge5,\\
 f_4h_{c_1}\prod_{i\ge2}h_{b_i}f_{c_i},&b_1=4.
 \end{cases}
\]
We prove $\E[p(T^L(\bar f_{a}P))^4]\le e^{-d/30}$ for every $L\ge1$.

\proofpart{The actual anchor occurs once}
Put $n=a+L$, $U=p\bar f_{n}^3$, $Z_j=A^jU$, $e_j=\E Z_j$ and
$W_{a,L}=\bar f_{a}Z_L$. Rowwise H\"older with density $\bar f_{a}$ gives
\begin{equation}\label{eq:compact-anchor-weight}
 (T^L(\bar f_{a}P))^4\le \bar f_{a+L}^3T^L(\bar f_{a}P^4),\qquad
 \E[p(T^L(\bar f_{a}P))^4]\le\E[W_{a,L}P^4].
\end{equation}
Its exact mass is $\E W_{a,L}=\E[p\bar f_{a+L}^4]\le1$ by
Section~\ref{loc:S20}. Also $e_0,e_1\le1$, $e_j\le2$ for $j\ge2$.
For $e_1$, use $pq\le1$ and the ordinary actual cubic bound
$\E \bar f_{n}^3\le1-101d/768$ proved in Section~\ref{q3:app:odd}.
It follows that $W_{a,L}\le4$ for $L=1,2$ and $W_{a,L}\le8$ otherwise.

For $d\le1/9$, the centered actual input satisfies
$\|U-\E U\|\le29\sqrt d/4$, as derived in
Section~\ref{loc:S28}. The centered $A$ recursion with
$B_0=(I-J)K(I-J)$ gives
\[
 Z_L-e_L=B_0^L(U-\E U)+\sum_{j=0}^{L-1}e_{L-1-j}B_0^jg.
\]
The finite coefficient bound from Section~\ref{loc:S27} yields centered
constants $457/80$, $569/128$, $12581/2560$ at $L=1,2,\ge3$.
Since $\|\bar f_{a}-\E \bar f_{a}\|\le3\sqrt d/2$, center the actual product by
comparison with the constant $(\E \bar f_{a})e_L$ to obtain constants
$517/40$, $761/64$, $16421/1280$, all below $13$.
Consequently
\begin{equation}\label{eq:compact-anchor-budget}
 \|W_{a,L}-\E W_{a,L}\|\le13\sqrt d.
\end{equation}
For an exponent $R$ with $\E R\le m_Rd$ and $\|R\|\le r_R\sqrt d$,
the scalar bound $e^{4R}\le1+4R+16cR^2$ therefore gives
\[
 \E[W_{a,L}e^{4R}]\le1+[4\max(m_R,0)+52r_R+128cr_R^2]d.
\]
The sharper squared norm may be retained in its final term.
The coefficients $c=2/3,13/18,4/5$ apply on
$|4R|\le3/4,1,5/4$, respectively; their series proofs are in
Sections~\ref{loc:S27}--\ref{loc:S30}.

\proofpart{First paired length at least five}
The actual product is $P\le e^{R-2td}$,
$R=\sum_iK^{b_i}G_{c_i-b_i}(K)g$.
For one pair $b_1=5$, the finite consecutive sum gives
$\E R\le0$, $\|R\|^2<d/64$ and $|4R|\le3/4$.
The cost is $47d/6$, leaving loss $d/6$ after $8d$.
For one pair $b_1\ge6$, use
$\|R\|\le21\sqrt d/250$, $\E R\le7d/100$;
the cost is at most $7/25+546/125+112896/187500<6$, times $d$.
For $t\ge2$, disjoint masks give
$\|R\|\le13\sqrt d/80$, $|R|\le\sqrt d/4$ and
$\E R\le21d/250$. With $c=2/3$ the cost is
$42/125+169/20+169/75=16559/1500<12$.
It leaves loss $[8(t-2)+4]d$.
These bounds apply for $d\le1/9$.
For $d\ge1/9$, the same row bound is $|R|\le3d/4$;
the exact mass $\E W_{a,L}\le1$ gives loss $(8t-3)d$ directly.

\proofpart{First paired length four}
Use the coefficients $C,D,E$ of Section~\ref{loc:S28}, with $c=c_1$.
Retaining the extra constant in $f_4=S_4-2d-dg$ gives
\[
 f_4h_c\le e^{R_4-4d},\qquad
 R_4=-\sum_{j=5}^cK^j\one+(C-d)g+Dr+Es.
\]
For $c=5$, $R_4=-K^4g+\delta_{2}g+\delta_{2}r$.
For $c\ge6$, its multiplier at $g$ is
\[
 Q_4=\alpha+(d-x)E-x(1-u)
                  +i\sqrt x[D-x(1-v)],\quad
 \alpha=d-\beta_c\in[0,\delta_{2}],\quad G_{c-2}=u+i\sqrt xv.
\]
Thus $\E R_4\le \delta_{3}$.
For $d\le1/9$ the full frequency bound is $|Q_4|^2\le9/200$,
and for $c=5$ it improves to $|Q_4|\le17/100$.
For the latter square $(\delta_{2}-x^2)^2+x\delta_{2}^2$ and use
$\theta^4+\theta^3/81<17/600$.
For $c\ge6$ set $D_{\max}=d(2\theta+d)<103/1000$.
On $x\le d$, the real modulus is below $3/50$ and the imaginary
coefficient below $103/1000$; on $d\le x\le1/3$, they are below
$3/20$ and $1/9$. The resulting square bounds are
$9/2500+(103/1000)^2/9$ and $9/400+1/243$, both below $9/200$.
On $1/3\le x\le\theta$ maximize the two convex quadratics in
$D\in[0,D_{\max}]$, using imaginary choices $x^2-D$ and $x^3$.
The common component bounds are
\[
 |\Re Q_4|\le x^2+(x-d)D,\qquad
 |\text{imaginary coefficient}|\le\max(x^2-D,x^3).
\]
Thus the two quadratics are $[x^2+(x-d)D]^2+x(x^2-D)^2$
and $[x^2+(x-d)D]^2+x^7$.
At $D=0$ the square is at most $\theta^4(1+\theta)<1/25$.
At $D=D_{\max}$ the first is at most $1/25+1/324$:
its excess over $x^4(1+x)$ is
$d^2k[-2x^2+k((x-d)^2+x)]\le1/324$, $k=2\theta+d<1$.
For the other choice the real modulus increases to below
$5927/30000<1/5$, and $\theta^7<1/500$ gives square below $21/500$.
This exhausts the frequency interval. Therefore
\[
 \|R_4\|^2\le9d/200,\quad\|R_4\|\le3\sqrt d/14,
 \quad |R_4|\le\|Q_4(K)\one\|\le3/14.
\]
For one pair use $c=13/18$ and retain the sharper squared norm.
The cost is $169/256+78/7+104/25=16-1657/44800$,
leaving loss $1657d/44800>d/30$.
For large $d$, $R_4=R_0+K^4\one-dg$ and
$|K^4\one-dg|\le\sqrt{\delta_{3}+d^2}<8d/5$ give
$|R_4|\le601d/160$; the loss is $39d/40$.

For at least two pairs put $R=R_4+S$, with later minimum at least six.
If $c_1=5$, the norm constants add as $17/100+21/200<9/32$;
if $c_1\ge6$, they add as $3/14+13/200<9/32$.
Also $\E R\le d/4$ and $|R|\le3/14+13/240<5/16$.
Use $c=4/5$ in the common budget: the cost is
$1+117/8+81/10=949/40$.
The actual loss $(8t+8)d$ leaves $[8(t-2)+11/40]d$.
For $d\ge1/9$, $|R|\le679d/160$ leaves
$[8(t-2)+281/40]d$.
Both parameter ranges include $1/9$; all margins imply the claimed
$d/30$ loss. The common rules restore equal pairs, $d=0$ and full reversal.
